% Source: gradient_descent.tex
\begin{notechapter}
% ---- chapter-local definitions ----
\definecolor{codegreen}{rgb}{0,0.6,0}
\definecolor{codegray}{rgb}{0.5,0.5,0.5}
\definecolor{codepurple}{rgb}{0.58,0,0.82}
\definecolor{backcolour}{rgb}{0.95,0.95,0.92}
\lstdefinestyle{mystyle}{
    backgroundcolor=\color{backcolour},
    commentstyle=\color{codegreen},
    keywordstyle=\color{magenta},
    numberstyle=\tiny\color{codegray},
    stringstyle=\color{codepurple},
    basicstyle=\ttfamily\footnotesize,
    breakatwhitespace=false,
    breaklines=true,
    captionpos=b,
    keepspaces=true,
    numbers=left,
    numbersep=5pt,
    showspaces=false,
    showstringspaces=false,
    showtabs=false,
    tabsize=2
}
\lstset{style=mystyle}
% -----------------------------------
\chapter{Gradient Descent and its Variants}\label{ch:optimizers}
\chaptersubtitle{Momentum, adaptive learning rates and decay}

\section{Introduction}
In machine learning, optimization techniques like momentum, 
adaptive learning rates, and learning rate decay are often 
combined to improve the performance of gradient-based 
algorithms. This document explains how these methods work 
individually, how they can be used together, and provides 
Python examples for practical implementation.

\section{Basic Gradient Descent}
Gradient descent is the foundation of many optimization 
algorithms. The goal is to minimize a loss function 
\( J(\theta) \) by iteratively updating the parameters 
\(\theta\) in the direction of the negative gradient of 
the loss function with respect to the parameters.

\subsection{Update Rule}
The basic update rule for gradient descent is:
\[
\theta_{t+1} = \theta_t - \eta \nabla_\theta J(\theta_t)
\]
where:
\begin{itemize}
    \item \(\theta_t\): Parameters at step \(t\)
    \item \(\eta\): Learning rate (step size)
    \item \(\nabla_\theta J(\theta_t)\): Gradient of the 
          loss function \(J\) with respect to \(\theta\) 
          at step \(t\)
\end{itemize}

\subsection{Intuition}
The learning rate \(\eta\) controls the size of the steps 
taken toward the minimum. A small learning rate leads to 
slow convergence, while a large learning rate can cause 
overshooting or divergence. The gradient 
\(\nabla_\theta J(\theta_t)\) points in the direction of 
the steepest ascent, so subtracting it moves the parameters 
in the direction of the steepest descent.

\subsection{Python Example: Basic Gradient Descent}
\begin{lstlisting}[language=Python]
import torch

# Define a simple model and loss function
model = torch.nn.Linear(1, 1)
criterion = torch.nn.MSELoss()
optimizer = torch.optim.SGD(model.parameters(), lr=0.01)

# Training loop
for epoch in range(100):
    optimizer.zero_grad()  # Clear gradients
    outputs = model(inputs)  # Forward pass
    loss = criterion(outputs, targets)  # Compute loss
    loss.backward()  # Backward pass (compute gradients)
    optimizer.step()  # Update parameters
\end{lstlisting}

\section{Curvature and the Hessian}
The curvature of a function is described by its second-order 
derivatives, which are captured in the Hessian matrix. The 
Hessian provides a way to measure the second-order behavior 
of a function, which is crucial for understanding how the 
function curves or bends in different directions.

\subsection{What is the Hessian?}
The Hessian matrix \( H \) of a function \( f(\mathbf{x}) \) 
is a square matrix of second-order partial derivatives. For 
a function \( f: \mathbb{R}^n \rightarrow \mathbb{R} \), the 
Hessian is defined as:

\[
H(f) = \begin{bmatrix}
\frac{\partial^2 f}{\partial x_1^2} & \frac{\partial^2 f}{\partial x_1 \partial x_2} & \cdots & \frac{\partial^2 f}{\partial x_1 \partial x_n} \\
\frac{\partial^2 f}{\partial x_2 \partial x_1} & \frac{\partial^2 f}{\partial x_2^2} & \cdots & \frac{\partial^2 f}{\partial x_2 \partial x_n} \\
\vdots & \vdots & \ddots & \vdots \\
\frac{\partial^2 f}{\partial x_n \partial x_1} & \frac{\partial^2 f}{\partial x_n \partial x_2} & \cdots & \frac{\partial^2 f}{\partial x_n^2}
\end{bmatrix}
\]

Each element \( H_{ij} \) of the Hessian represents the 
second-order partial derivative of \( f \) with respect to 
\( x_i \) and \( x_j \).

\subsection{How Does the Hessian Relate to Curvature?}
The Hessian describes the local curvature of the function 
\( f \) at a given point \( \mathbf{x} \). Specifically:
\begin{itemize}
    \item \textbf{Positive Definite Hessian}: If the Hessian 
          is positive definite at a point \( \mathbf{x} \), 
          the function curves upward in all directions, and 
          \( \mathbf{x} \) is a local minimum.
    \item \textbf{Negative Definite Hessian}: If the Hessian 
          is negative definite at a point \( \mathbf{x} \), 
          the function curves downward in all directions, and 
          \( \mathbf{x} \) is a local maximum.
    \item \textbf{Indefinite Hessian}: If the Hessian is 
          indefinite (i.e., it has both positive and negative 
          eigenvalues), the function has a saddle point at 
          \( \mathbf{x} \), meaning it curves upward in some 
          directions and downward in others.
\end{itemize}

\subsection{Eigenvalues of the Hessian and Curvature}
The eigenvalues of the Hessian matrix provide quantitative 
information about the curvature:
\begin{itemize}
    \item \textbf{Large Positive Eigenvalues}: Indicate strong 
          upward curvature in the corresponding eigenvector 
          directions.
    \item \textbf{Small Positive Eigenvalues}: Indicate weak 
          upward curvature.
    \item \textbf{Negative Eigenvalues}: Indicate downward 
          curvature (concavity) in the corresponding eigenvector 
          directions.
    \item \textbf{Zero Eigenvalues}: Indicate no curvature 
          (flatness) in the corresponding eigenvector directions.
\end{itemize}

\subsection{Practical Use of the Hessian in Optimization}
The Hessian is used in various ways in optimization:
\begin{itemize}
    \item \textbf{Second-Order Optimization Methods}: Methods 
          like Newton's method use the Hessian to take into 
          account the curvature of the loss function. The update 
          rule for Newton's method is:
          \[
          \mathbf{x}_{t+1} = \mathbf{x}_t - H^{-1} \nabla f(\mathbf{x}_t),
          \]
          where \( H^{-1} \) is the inverse of the Hessian matrix. 
          These methods converge faster than first-order methods 
          (like gradient descent) because they account for 
          curvature, but they are computationally expensive due 
          to the need to compute and invert the Hessian.

    \item \textbf{Saddle Points and Ill-Conditioning}: In 
          high-dimensional optimization problems, saddle points 
          (where the Hessian has both positive and negative 
          eigenvalues) are common. Understanding the Hessian 
          helps in designing algorithms to escape saddle points. 
          The condition number of the Hessian (the ratio of the 
          largest to smallest eigenvalue) determines how 
          "ill-conditioned" the optimization problem is. A high 
          condition number indicates that the curvature varies 
          significantly across different directions, making 
          optimization more challenging.

    \item \textbf{Approximations of the Hessian}: In large-scale 
          problems, computing the full Hessian is often impractical. 
          Instead, approximations like the BFGS algorithm or 
          L-BFGS are used to estimate the Hessian or its inverse.
\end{itemize}

\subsection{Example: Hessian in Machine Learning}
In machine learning, the Hessian is often used to analyze the 
loss surface of a model. For example:
\begin{itemize}
    \item \textbf{Neural Networks}: The Hessian of the loss 
          function with respect to the model parameters can 
          reveal insights into the geometry of the loss surface, 
          such as the presence of flat regions (small eigenvalues) 
          or sharp minima (large eigenvalues).
    \item \textbf{Regularization}: The Hessian can be used to 
          analyze the effect of regularization terms (e.g., L2 
          regularization) on the curvature of the loss function.
\end{itemize}

\subsection{BFGS (Broyden–Fletcher–Goldfarb–Shanno)}
BFGS is a quasi-Newton optimization algorithm that approximates the Hessian matrix 
using gradient information from previous iterations. It is widely used for solving 
unconstrained nonlinear optimization problems.

\subsubsection{Key Features}
\begin{itemize}
    \item \textbf{Quasi-Newton Method}: BFGS approximates the Hessian matrix (or its 
          inverse) without explicitly computing it, making it efficient for 
          high-dimensional problems.
    \item \textbf{Hessian Approximation}: The Hessian approximation is updated 
          iteratively using the secant equation, ensuring that it satisfies curvature 
          conditions.
    \item \textbf{Positive Definiteness}: BFGS ensures that the Hessian approximation 
          remains positive definite, guaranteeing a descent direction.
    \item \textbf{Superlinear Convergence}: BFGS exhibits superlinear convergence 
          under certain conditions.
\end{itemize}

\subsubsection{BFGS Update Formula}
The Hessian approximation \( B_k \) is updated as:
\[
B_{k+1} = B_k + \frac{y_k y_k^T}{y_k^T s_k} - \frac{B_k s_k s_k^T B_k}{s_k^T B_k s_k},
\]
where:
\begin{itemize}
    \item \( s_k = x_{k+1} - x_k \): The change in the parameter vector.
    \item \( y_k = \nabla f(x_{k+1}) - \nabla f(x_k) \): The change in the gradient.
\end{itemize}

\subsubsection{Variants of BFGS}
\begin{itemize}
    \item \textbf{L-BFGS}: A memory-efficient variant that stores only a few vectors 
          to approximate the Hessian, making it suitable for large-scale problems.
    \item \textbf{BFGS with Trust Regions}: Incorporates trust-region methods to 
          improve robustness and convergence.
\end{itemize}

\subsubsection{Practical Use}
BFGS is widely used in machine learning and scientific computing. For example:
\begin{itemize}
    \item \textbf{SciPy}: The `scipy.optimize.minimize` function provides an 
          implementation of BFGS.
    \item \textbf{PyTorch}: The `torch.optim.LBFGS` optimizer implements L-BFGS for 
          training neural networks.
\end{itemize}

\subsubsection{Python Example: BFGS in SciPy}
\begin{lstlisting}[language=Python]
from scipy.optimize import minimize

# Define the objective function
def objective(x):
    return x[0]**2 + x[1]**2 + x[0]*x[1] + x[0] + x[1]

# Initial guess
x0 = [0.0, 0.0]

# Minimize using BFGS
result = minimize(objective, x0, method='BFGS')

print("Optimal solution:", result.x)
print("Objective value at solution:", result.fun)
\end{lstlisting}

\section{Momentum}
Momentum is a technique that helps accelerate gradient descent
by adding a fraction of the previous update vector to the current
update. This helps in smoothing the path taken by the optimizer
and can lead to faster convergence, especially in the presence
of noisy gradients or when the loss surface has high curvature.

\subsection{Standard Momentum}
The update rule for standard momentum is:
\[
v_t = \gamma v_{t-1} + \eta \nabla_\theta J(\theta)
\]
\[
\theta = \theta - v_t
\]
where:
\begin{itemize}
    \item \(v_t\): Update vector at step \(t\)
    \item \(\gamma\): Momentum term (e.g., 0.9)
    \item \(\eta\): Learning rate
    \item \(\nabla_\theta J(\theta)\): Gradient of the loss \(J\)
\end{itemize}

\subsection{Nesterov Momentum}
Nesterov momentum improves standard momentum by evaluating 
the gradient at the look-ahead position:
\[
v_t = \gamma v_{t-1} + \eta \nabla_\theta J(\theta - \gamma v_{t-1})
\]
\[
\theta = \theta - v_t
\]
This correction helps reduce oscillations and improves 
convergence.

\subsection{Python Example: SGD with Momentum}
\begin{lstlisting}[language=Python]
import torch
optimizer = torch.optim.SGD(model.parameters(), 
                            lr=0.1, momentum=0.9)
\end{lstlisting}

\section{Adaptive Learning Rates}
Adaptive learning rate methods, such as AdaGrad, RMSProp,
and Adam, adjust the learning rate for each parameter based
on the history of gradients. This allows the optimizer to take
larger steps in directions where the gradient is consistently
small and smaller steps in directions where the gradient is
large or fluctuating.

\subsection{AdaGrad}
AdaGrad adapts the learning rate by dividing it by the 
square root of the sum of squared gradients:
\[
g_{t,i} = \nabla_\theta J(\theta_i)
\]
\[
G_{t,i} = G_{t-1,i} + g_{t,i}^2
\]
\[
\theta_{t+1,i} = \theta_{t,i} - \frac{\eta}{\sqrt{G_{t,i} + \epsilon}} g_{t,i}
\]
where \(G_{t,i}\) accumulates squared gradients and 
\(\epsilon\) is a small constant for numerical stability.

\subsection{RMSProp}
RMSProp modifies AdaGrad by using an exponential moving 
average of squared gradients:
\[
G_{t,i} = \gamma G_{t-1,i} + (1 - \gamma) g_{t,i}^2
\]
\[
\theta_{t+1,i} = \theta_{t,i} - \frac{\eta}{\sqrt{G_{t,i} + \epsilon}} g_{t,i}
\]
This prevents the learning rate from shrinking too much.

\subsection{Adam}
Adam (Adaptive Moment Estimation) is a popular
optimization algorithm that combines the benefits of both
momentum and adaptive learning rates. It uses exponential
moving averages of both the gradients (momentum) and the
squared gradients (adaptive learning rate) to update the
parameters.

\[
m_t = \beta_1 m_{t-1} + (1 - \beta_1) g_t
\]
\[
v_t = \beta_2 v_{t-1} + (1 - \beta_2) g_t^2
\]
\[
\hat{m}_t = \frac{m_t}{1 - \beta_1^t}, \quad \hat{v}_t = \frac{v_t}{1 - \beta_2^t}
\]
\[
\theta_{t+1} = \theta_t - \frac{\eta}{\sqrt{\hat{v}_t} + \epsilon} \hat{m}_t
\]
Adam is robust and widely used in practice.

\subsection{Python Example: Adam Optimizer}
\begin{lstlisting}[language=Python]
import torch
optimizer = torch.optim.Adam(model.parameters(), 
                             lr=0.001)
\end{lstlisting}

\section{Learning Rate Decay}
Learning rate decay reduces the learning rate over time to 
fine-tune the model and improve convergence. This technique 
is particularly useful in the later stages of training, 
where smaller updates can help the model converge to a 
better minimum.

\subsection{Why Use Learning Rate Decay?}
Learning rate decay gradually reduces the learning rate over 
time, which can help:
\begin{itemize}
    \item \textbf{Fine-tune the model}: A smaller learning rate 
          in later stages allows the optimizer to make smaller, 
          more precise updates, helping the model converge to a 
          better minimum.
    \item \textbf{Prevent overshooting}: As the model gets closer 
          to the optimal solution, large updates might cause the 
          optimizer to overshoot the minimum. Decaying the learning 
          rate mitigates this risk.
    \item \textbf{Improve stability}: In adaptive methods like Adam, 
          the learning rate is adjusted per parameter, but a global 
          decay can still help stabilize training over long periods.
\end{itemize}

\subsection{Step Decay}
Step decay reduces the learning rate by a factor every few 
epochs:
\[
\eta_t = \eta_0 \times \gamma^{\lfloor t / s \rfloor}
\]
where \(s\) is the step size and \(\gamma\) is the decay rate.

\subsection{Exponential Decay}
Exponential decay reduces the learning rate continuously:
\[
\eta_t = \eta_0 \times e^{-kt}
\]
where \(k\) controls the decay rate.

\subsection{Cosine Decay}
Cosine decay reduces the learning rate following a cosine 
curve:
\[
\eta_t = \eta_0 \times \frac{1 + \cos(\pi t / T)}{2}
\]
where \(T\) is the total number of steps.

\subsection{Python Example: Learning Rate Scheduler}
\begin{lstlisting}[language=Python]
import torch
optimizer = torch.optim.SGD(model.parameters(), 
                            lr=0.1, momentum=0.9)
scheduler = torch.optim.lr_scheduler.StepLR(
    optimizer, step_size=30, gamma=0.1
)
\end{lstlisting}

\section{Combining Techniques}
Momentum, adaptive learning rates, and learning rate decay 
can be combined to enhance optimization. Below, we discuss 
how to effectively combine these methods.

\subsection{Momentum with Decay}
When using momentum, apply a learning rate schedule to the 
base learning rate:
\[
v_t = \gamma v_{t-1} + \eta_t \nabla_\theta J(\theta)
\]
\[
\theta = \theta - v_t
\]
where \(\eta_t\) is the decayed learning rate.

\subsection{Adaptive Methods with Decay}
For adaptive methods like Adam, apply a global learning 
rate decay:
\[
\eta_t = \eta_0 \times \text{decay}(t)
\]
\[
\theta_{t+1} = \theta_t - \frac{\eta_t}{\sqrt{\hat{v}_t} + \epsilon} \hat{m}_t
\]
This combines the benefits of adaptive learning rates and 
decay.

\subsection{Python Example: Adam with Learning Rate Decay}
\begin{lstlisting}[language=Python]
import torch
optimizer = torch.optim.Adam(model.parameters(), 
                             lr=0.001)
scheduler = torch.optim.lr_scheduler.ExponentialLR(
    optimizer, gamma=0.95
)
\end{lstlisting}

\subsection{Practical Considerations}
When combining momentum, adaptive learning rates, and decay, 
consider the following:
\begin{itemize}
    \item \textbf{Initial Learning Rate}: When using decay, the 
          initial learning rate should be set higher to account 
          for the reduction over time.
    \item \textbf{Decay Schedule}: The choice of decay schedule 
          (e.g., step, exponential, cosine) depends on the problem 
          and dataset. Experimentation is often required to find 
          the best schedule.
    \item \textbf{Adaptive Methods}: While adaptive methods like 
          Adam are less sensitive to the initial learning rate, 
          adding decay can still help in fine-tuning the model.
\end{itemize}

\section{Conclusion}
Momentum, adaptive learning rates, and learning rate decay 
are complementary techniques. By combining them, you can 
achieve faster convergence, better stability, and improved 
performance in machine learning optimization. The provided 
Python examples demonstrate how to implement these methods 
in practice.
\end{notechapter}
